\label{sec:referencial}

Este capítulo reúne a base conceitual e empírica da proposta. Os temas foram organizados nos seis eixos usados na revisão narrativa descrita no Capítulo~\ref{sec:metodologia}: superlotação dos serviços de urgência, perfil dos pacientes de baixo risco, Protocolo de Manchester, quiosques e autoatendimento, telemedicina e seu marco regulatório, e interoperabilidade e proteção de dados. Ao final, apresentam-se os fundamentos de arquitetura de software usados na modelagem e os trabalhos correlatos.

\subsecao{Superlotação dos serviços de urgência}

A superlotação de um serviço de emergência ocorre quando a demanda supera a capacidade de atender com qualidade e em tempo adequado. Savioli et al. \cite{savioli2022} organizam suas causas no modelo de entrada, processamento e saída (\textit{input-throughput-output}). Os fatores de entrada determinam o acesso do paciente ao serviço, incluindo o volume de chegadas e a gravidade e complexidade dos casos. Os fatores de processamento correspondem ao tempo entre o início do cuidado e a definição da conduta. Os fatores de saída incluem a permanência de pacientes já atendidos na emergência à espera de leito (\textit{boarding}), a disponibilidade de leitos hospitalares e a capacidade de alta ou transferência. Os autores destacam o \textit{boarding} entre os fatores causais mais importantes, por consumir recursos da emergência e perpetuar a superlotação \cite{savioli2022}.

Duas revisões de revisões ampliam esse quadro. A primeira reuniu medidas, causas e danos associados à superlotação \cite{pearce2023}. A segunda sintetizou 27 revisões sistemáticas, que somavam 308 estudos primários, sobre intervenções organizadas pelos mesmos três níveis do modelo. As intervenções de entrada e de saída apresentaram os melhores resultados nas métricas de superlotação, e a grande variação dos resultados entre contextos levou os autores a recomendar que os fatores locais sejam avaliados antes da adoção de qualquer solução \cite{pearce2024}.

No Brasil, Freitas et al. \cite{freitas2024} descrevem a superlotação das unidades de emergência como uma realidade com tendência de piora e relacionam suas implicações à segurança do paciente e ao trabalho cotidiano da equipe de saúde. Para este trabalho, esse conjunto de evidências tem duas consequências. A primeira é que uma intervenção sobre o fluxo de pacientes de baixo risco atua sobre a entrada e o processamento, e não sobre a saída, que é onde está parte relevante do problema. A segunda é que qualquer proposta precisa ser avaliada no contexto local antes de se afirmar algum efeito sobre a superlotação. As duas consequências limitam o alcance das conclusões deste trabalho e são retomadas na discussão.

\subsecao{Baixo risco e Protocolo de Manchester}

O Protocolo de Manchester (\textit{Manchester Triage System}, MTS) classifica o paciente em cinco níveis de urgência, identificados por cores, cada um associado a um tempo máximo até o primeiro atendimento médico: vermelho, imediato; laranja, 10 minutos; amarelo, 60 minutos; verde, 2 horas; e azul, 4 horas \cite{stormversloot2014}. A classificação é feita por profissional treinado, a partir de fluxogramas e discriminadores clínicos. Neste trabalho, os pacientes de baixo risco são os classificados como verde (pouco urgente) ou azul (não urgente).

Storm-Versloot et al. \cite{stormversloot2014} compararam, em estudo do tipo antes e depois em um hospital universitário holandês, o tempo de espera, o tempo de tratamento, o tempo de permanência e a satisfação dos pacientes antes e depois da implantação do MTS. Com 1.808 pacientes em duas séries consecutivas, o tempo de espera não diminuiu, embora tenha passado a ser mais bem distribuído entre os níveis de urgência, e o tempo de tratamento aumentou de forma significativa. Na satisfação, a avaliação do tempo de espera e da sensação de resolução do problema melhorou, mas a da oferta de informações e da oportunidade de explicar a queixa piorou. Os autores concluíram que implantar o MTS, sozinho, não basta para melhorar a eficiência e a qualidade do serviço. A leitura relevante para esta proposta é que a triagem, sozinha, organiza a prioridade, mas não cria capacidade de atendimento.

Sorz-Nechay et al. \cite{sorznechay2025} avaliaram um processo de encaminhamento de pacientes pediátricos de baixo risco em Viena. Em 2015, o departamento de emergência pediátrica atingiu sua capacidade com cerca de 25 mil atendimentos, dos quais 95\% eram de procura espontânea e 75\% classificados nos níveis verde e azul. A partir de 2016, esses pacientes passaram a ser encaminhados a consultórios pediátricos e a serviços secundários, e o volume caiu para 16.500 atendimentos em 2017. Entre junho de 2017 e maio de 2018 foram incluídos 2.394 pacientes encaminhados, dos quais 568 responderam ao questionário de seguimento. Seis deles (1,1\%) precisaram de internação em até duas semanas, nenhum em terapia intensiva, e a satisfação das famílias foi alta.

Dois aspectos do desenho desse estudo importam para a arquitetura proposta. Primeiro, a cor atribuída pelo MTS não bastava para o encaminhamento: a decisão era tomada pelo médico do serviço a partir da documentação da triagem, e havia critérios de exclusão, como recém-nascidos, certas doenças crônicas, suspeita de doença altamente infecciosa, espera superior a 15 minutos pela decisão médica e barreiras linguísticas importantes \cite{sorznechay2025}. Segundo, o desfecho de segurança foi medido por um critério objetivo, a internação em 14 dias. Ambos os aspectos sustentam a decisão de tratar a classificação verde ou azul como condição necessária, mas não suficiente, para o encaminhamento ao canal remoto.

Quanto às razões da procura, Schmiedhofer et al. \cite{schmiedhofer2016} entrevistaram pacientes de baixa acuidade em serviços de emergência alemães, urbanos e rurais. Os motivos predominantes foram a conveniência, ligada à disponibilidade em tempo integral e à concentração de recursos diagnósticos em um só lugar, e a ansiedade em relação à própria saúde, além de restrições de tempo ligadas ao trabalho. Esses achados indicam que um canal alternativo precisa ser oferecido no momento e no local em que o paciente já está, e precisa dar a ele a percepção de ter sido avaliado por um médico.

\subsecao{Telemedicina: expansão e marco regulatório}

A telemedicina é o exercício da medicina mediado por tecnologias digitais, de informação e de comunicação \cite{cfm2314}. Em revisão de escopo sobre tecnologias digitais na atenção primária durante a pandemia, Ndayishimiye, Lopes e Middleton \cite{ndayishimiye2023} identificaram 19 tipos de tecnologias, agrupados em cinco funções: prestação geral de serviços (teleconsulta, diagnóstico e prescrição eletrônicos), promoção de comportamentos e letramento em saúde digital, vigilância, vacinação e medicamentos, e apoio à decisão nos sistemas de saúde. Os autores apontam como necessidades a padronização técnica, a garantia de qualidade e modelos de financiamento.

No Brasil, a Lei n.º 13.989/2020 autorizou o uso da telemedicina durante a crise causada pelo coronavírus \cite{lei13989}. Donida, Costa e Scherer \cite{donida2021} relatam que, em menos de quatro meses, o serviço TeleSUS realizou 7,4 milhões de atendimentos remotos, e que a estratégia Conecte SUS e a RNDS foram redirecionadas para apoiar o enfrentamento da pandemia.

O marco atual tem três normas principais. A Resolução CFM n.º 2.314/2022 \cite{cfm2314} define as modalidades de telemedicina, entre elas a teleconsulta, e estabelece, entre outros pontos:
\begin{itemize}
    \item o registro do atendimento em prontuário, físico ou em sistema de registro eletrônico de saúde que atenda ao nível de garantia de segurança 2 (NGS2) e ao padrão ICP-Brasil, com preservação da confidencialidade, integridade e privacidade dos dados (art. 3.º);
    \item a autonomia do médico para decidir sobre o uso da telemedicina e indicar o atendimento presencial, com uma linha de cuidado que preveja o atendimento presencial nos casos de risco (art. 4.º);
    \item a consulta presencial como padrão de referência, a obrigação de informar as limitações do atendimento remoto e a possibilidade de interrupção por iniciativa do médico ou do paciente (art. 6.º);
    \item a emissão de documentos, como prescrições e atestados, com identificação do médico e do paciente, data e hora, assinatura digital ICP-Brasil e indicação de que foram emitidos por telemedicina (art. 13);
    \item o consentimento livre e esclarecido do paciente para o atendimento e para a transmissão de imagens e dados, com termo arquivado no prontuário (art. 15);
    \item a vedação de que os serviços a distância substituam a assistência presencial garantida pelos princípios do SUS (art. 19).
\end{itemize}

A Lei n.º 14.510/2022 incluiu na Lei n.º 8.080/1990 os artigos 26-A a 26-H, que disciplinam a telessaúde em todo o território nacional \cite{lei14510}. Entre os princípios do art. 26-A estão a autonomia do profissional, o consentimento livre e informado do paciente, o direito de recusa ao atendimento a distância com garantia do atendimento presencial sempre que solicitado, e a confidencialidade dos dados. O art. 26-C assegura ao profissional a liberdade de decidir sobre o uso da telessaúde e de indicar o atendimento presencial. O art. 26-G exige o consentimento livre e esclarecido do paciente ou de seu representante legal, sob responsabilidade do profissional, e a observância de outras normas, como a LGPD.

Por fim, a Portaria GM/MS n.º 3.691/2024 instituiu a Ação Estratégica SUS Digital - Telessaúde, com modalidades como teleconsulta, teletriagem e teleinterconsulta, e definiu como ação estruturante que as plataformas de telessaúde sejam interoperáveis com a RNDS, com os prontuários eletrônicos e com a plataforma Meu SUS Digital, seguindo os padrões de interoperabilidade definidos pelo Ministério da Saúde \cite{portaria3691}.

\subsecao{Quiosques de autoatendimento e autotriagem}

Quiosques de autoatendimento são dispositivos em que o próprio usuário registra informações por tela sensível ao toque, sem mediação direta de um profissional. Em serviços de emergência, são usados para o registro da chegada (\textit{self-check-in}) e para a coleta de informações clínicas destinadas à pré-classificação (\textit{e-triage}) \cite{joseph2023}.

Joseph et al. \cite{joseph2023} revisaram nove estudos sobre quiosques em serviços de emergência. Seis tiveram resultados positivos, como a redução do tempo antes da triagem, e três tiveram resultados negativos, como baixa adesão. Em um dos estudos incluídos, com dados nacionais dos Estados Unidos, serviços com quiosques foram associados a tempos de espera quase 57\% menores. Em outro, pacientes atendidos em serviços com dispositivos de autoatendimento tiveram menor chance de espera prolongada (razão de chances 0,56; IC 95\% 0,41--0,83) e de \textit{boarding} prolongado (razão de chances 0,55; IC 95\% 0,35--0,85). As barreiras descritas pelos autores são especialmente relevantes para o projeto: critérios de elegibilidade muito restritivos, exigência de algum letramento em saúde para responder ao questionário, barreiras de idioma, resistência da enfermagem em encaminhar pacientes e desconfiança dos médicos em relação às avaliações feitas por algoritmo. Os autores recomendam os quiosques como complemento, e não substituto, da triagem conduzida pela enfermagem.

A limitação da autotriagem foi medida diretamente por Dickson et al. \cite{dickson2022}, em estudo retrospectivo de oito meses em dois serviços de emergência do Reino Unido, com 25.333 pares de classificações de adultos ambulatoriais. A concordância entre a autotriagem eletrônica e a triagem de enfermagem pelo MTS foi baixa (kappa ponderado de 0,14). A autotriagem classificou como mais urgente do que a enfermagem em 59,2\% dos casos e como menos urgente em 10,1\%. Sua sensibilidade para desfechos de alta acuidade foi maior (88,5\% contra 53,8\%), assim como a especificidade para pacientes de baixo risco (88,5\% contra 80,6\%). Os autores não propõem a substituição da triagem de enfermagem. Para esta proposta, a baixa concordância é o argumento central para que o totem não classifique risco: ele recebe o paciente depois da classificação feita pela enfermagem.

\subsecao{Proteção de dados pessoais em saúde}

A LGPD classifica o dado referente à saúde como dado pessoal sensível (art. 5.º, II) e estabelece hipóteses específicas para o seu tratamento (art. 11) \cite{lgpd2018}. Uma delas é a tutela da saúde, exclusivamente em procedimento realizado por profissionais de saúde, serviços de saúde ou autoridade sanitária (art. 11, II, \textit{f}). Outra é o consentimento específico e destacado do titular (art. 11, I). O tratamento também se submete aos princípios do art. 6.º, entre eles finalidade, necessidade, transparência, segurança, prevenção e responsabilização, e às obrigações de registro das operações (art. 37), de elaboração de relatório de impacto quando solicitado (art. 38), de indicação de encarregado (art. 41), de adoção de medidas de segurança (art. 46) e de comunicação de incidentes à autoridade nacional e ao titular (art. 48).

Leme e Blank \cite{leme2020} analisam a LGPD na área da saúde e mostram que o uso de meios digitais expõe o titular a riscos de negligência no tratamento, de comercialização indevida e de vazamento. Os autores defendem o consentimento qualificado como elemento essencial e a adoção de sistemas seguros. Dourado e Aith \cite{dourado2022} destacam o direito à revisão de decisões tomadas unicamente com base em tratamento automatizado (art. 20) como marco inicial da regulação da inteligência artificial em saúde no Brasil.

Para o desenho de um sistema assistencial público, uma distinção é decisiva e costuma ser confundida: a base legal do tratamento de dados e o consentimento para o atendimento são institutos diferentes. Em um atendimento prestado por serviço de saúde, a base legal adequada é a tutela da saúde, e não o consentimento. Já o consentimento exigido pela Resolução CFM n.º 2.314/2022 e pela Lei n.º 14.510/2022 é um ato clínico de autorização do atendimento a distância \cite{cfm2314, lei14510}. Por isso, a recusa ou a revogação desse consentimento interrompe a teleconsulta, mas não elimina o dever de registrar e guardar o prontuário. A Lei n.º 13.787/2018 disciplina a guarda do prontuário em meio eletrônico e fixa prazo mínimo de 20 anos a partir do último registro \cite{lei13787}, e a própria LGPD admite a conservação dos dados para cumprimento de obrigação legal (art. 16, I) \cite{lgpd2018}.

\subsecao{Ecossistema digital do SUS}

\subsubsecao{e-SUS Atenção Primária e Prontuário Eletrônico do Cidadão}

A estratégia e-SUS Atenção Primária (e-SUS APS) informatiza o registro da atenção básica, e o Prontuário Eletrônico do Cidadão (PEC) é seu sistema mais abrangente. Celuppi et al. \cite{celuppi2024} descrevem dez anos de evolução do PEC: o número de unidades básicas que o utilizavam passou de 8.930, em 2017, para 26.091, em 2022, quando o PEC respondia por 59,93\% dos registros eletrônicos da atenção primária. Os autores registram que o sistema não foi concebido originalmente para integração com sistemas de terceiros e que a integração com a atenção especializada e hospitalar continua sendo uma lacuna, e propõem ampliar o compartilhamento de dados pela RNDS.

Coelho Neto, Andreazza e Chioro \cite{neto2021} identificaram 31 sistemas nacionais de informação usados na atenção primária: 12 estavam integrados ao PEC, 4 parcialmente integrados e 15 sem interfaces unificadas, uma fragmentação que os autores associam à separação entre assistência e vigilância no Ministério da Saúde. Cielo et al. \cite{cielo2022} analisaram dados oficiais de 2013 a 2019 e classificaram a implantação da estratégia e-SUS AB em 2019 como completa em 20,2\% dos municípios, parcial em 32,9\%, inicial em 39,1\% e inexistente em 7,8\%. Esses dados mostram que a integração com o registro da atenção primária não pode ser pressuposta: ela deve ser tratada como opcional e dependente da realidade de cada unidade.

\subsubsecao{Rede Nacional de Dados em Saúde}

A Portaria GM/MS n.º 1.434/2020 instituiu o Programa Conecte SUS e definiu a RNDS como plataforma nacional de integração e interoperabilidade de informações em saúde \cite{portaria1434}. A portaria distingue interoperabilidade semântica, baseada em modelos de informação e vocabulários padronizados, e interoperabilidade sintática, baseada em modelos computacionais para troca de informações sem perda de significado, e dá preferência a padrões abertos. Donida, Costa e Scherer \cite{donida2021} descrevem a RNDS como uma camada federada de interoperabilidade cuja primeira fase, durante a pandemia, compartilhou resultados de exames de COVID-19 com base no HL7 FHIR.

A especificação técnica da rede é publicada em duas frentes. O Guia de Implementação da RNDS reúne os perfis FHIR nacionais, como BRIndividuo, BRProfissional, BRLotacaoProfissional e BREstabelecimentoSaude, além dos sistemas de nomenclatura para CNS, CPF, CNES, CBO, CID-10 e CIAP-2, sobre a versão 4.0.1 do FHIR \cite{rndsig}. O Guia de Integração descreve o funcionamento prático: o sistema que se integra, chamado de conector, obtém um token de acesso no serviço de autenticação usando o certificado digital do estabelecimento de saúde, reutiliza o token durante sua validade de 30 minutos e o envia, nas requisições aos serviços clínicos, no cabeçalho \texttt{X-Authorization-Server}, enquanto o cabeçalho \texttt{Authorization} leva o CNS do profissional em nome de quem a requisição é feita \cite{rndsconector, rndserros}. Há um ambiente de homologação único e, em produção, um endereço de serviços clínicos por unidade federativa \cite{rndsambientes}. Os erros são devolvidos como recursos \texttt{OperationOutcome}, com mensagens de negócio codificadas \cite{rndserros}.

O Registro de Atendimento Clínico (RAC) foi instituído como modelo de informação pela Portaria SAES/MS n.º 234/2022, de adoção obrigatória nas esferas pública, suplementar e privada \cite{portaria234}. O modelo define as seções do registro e sua obrigatoriedade, como identificação do indivíduo por CNS ou CPF, caracterização do atendimento, profissionais, motivo do atendimento, sinais vitais, problemas e diagnósticos avaliados codificados em CID ou CIAP, procedimentos, prescrição, plano de cuidados e desfecho. Em 2025, o Ministério da Saúde apresentou uma especialização do RAC para teleconsulta, com a inclusão da modalidade de telessaúde, cujo modelo computacional estava previsto para desenvolvimento e piloto no fim daquele ano \cite{ms2025modelos}. Essa informação é relevante porque indica que parte da especificação necessária para este caso de uso ainda estava em construção na data das fontes consultadas.

\subsecao{Arquitetura de software para telemedicina}

\subsubsecao{Conceitos usados na modelagem}

A arquitetura de um sistema de software descreve seus elementos, as relações entre eles e as decisões que orientam seu desenho. As qualidades que a arquitetura deve atender foram organizadas, neste trabalho, segundo o modelo de qualidade de produto da norma ISO/IEC 25010, que inclui adequação funcional, eficiência de desempenho, compatibilidade, usabilidade, confiabilidade, segurança, manutenibilidade e portabilidade \cite{iso25010}.

Para a documentação foi usado o modelo C4, que descreve a arquitetura em níveis de abstração sucessivos: contexto, com o sistema e seus usuários e sistemas externos; contêineres, com as aplicações e armazenamentos que compõem o sistema; componentes, com a estrutura interna de um contêiner; e código \cite{brownc4}. A avaliação por cenários adotada no Capítulo~\ref{sec:metodologia} inspira-se no método ATAM (\textit{Architecture Tradeoff Analysis Method}), que confronta a arquitetura com cenários que expressam atributos de qualidade e examina a resposta prevista, os riscos e os compromissos entre atributos \cite{kazman1998}.

\subsubsecao{HL7 FHIR}

O HL7 FHIR (\textit{Fast Healthcare Interoperability Resources}) organiza a informação em saúde em recursos, como paciente, profissional, encontro, observação e condição, que podem ser trocados por interfaces REST e representados em JSON ou XML \cite{hl7fhirr4}. Os recursos podem ser agrupados em um \textit{Bundle}, inclusive do tipo documento, e restringidos por perfis que fixam elementos obrigatórios, cardinalidades e terminologias para um contexto específico. Adotar o FHIR, portanto, não basta para garantir interoperabilidade com a RNDS: é preciso seguir os perfis, terminologias e regras de envio publicados para a rede.

\subsubsecao{Tendências e lacunas}

Jat et al. \cite{jat2024} revisaram 17 artigos sobre arquitetura de software em telemedicina e identificaram como tendências a computação em nuvem, sistemas descentralizados, \textit{blockchain} e HL7 FHIR. Segurança, interoperabilidade e escalabilidade foram as preocupações recorrentes. A revisão, porém, trata a conformidade regulatória de forma superficial, e muitos dos estudos incluídos são propostas conceituais ou protótipos, com pouca avaliação junto a pacientes.

Gøeg et al. \cite{goeg2018} propuseram uma arquitetura de telemedicina baseada em três princípios: um núcleo com módulos, inclusive para banco de dados e segurança; baixo acoplamento entre módulos por inversão de controle; e HL7 FHIR como padrão principal, por se basear em tecnologias web. A viabilidade foi avaliada com uma implementação de código aberto, escrita em TypeScript sobre o \textit{framework} Express e o FHIR DSTU2, com testes unitários dos módulos. Os testes de ponta a ponta ficaram pendentes da primeira implementação clínica. A ideia de isolar a integração em um módulo próprio é adotada nesta proposta no gateway FHIR.

\subsecao{Escores de alerta precoce}

O \textit{National Early Warning Score 2} (NEWS2), do Royal College of Physicians, padroniza a avaliação da gravidade de doenças agudas a partir de seis parâmetros fisiológicos: frequência respiratória, saturação de oxigênio, pressão arterial sistólica, frequência de pulso, nível de consciência ou confusão nova e temperatura, com peso adicional para uso de oxigênio suplementar \cite{rcp2017news2}. A pontuação agregada define a resposta clínica: de 1 a 4, avaliação por enfermeiro competente; 3 em um único parâmetro, revisão urgente por clínico; de 5 a 6, revisão urgente e frequência mínima de monitoramento por hora; 7 ou mais, avaliação de emergência e monitoramento contínuo. O NEWS2 foi concebido para o sistema de saúde britânico, e não especificamente para pacientes de baixo risco em espera. Por isso, neste trabalho ele é usado apenas como referência para gatilhos de deterioração, que precisam de validação clínica antes de qualquer uso.

\subsecao{Trabalhos correlatos}

A Tabela~\ref{tab:correlatos} compara os trabalhos mais próximos da proposta nas dimensões que ela pretende articular.

\begin{table}[H]
\centering
\caption{Comparação entre trabalhos correlatos e a proposta}
\label{tab:correlatos}
\footnotesize
\begin{tabularx}{\textwidth}{L{2.6cm} Y Y L{2.7cm}}
\toprule
\textbf{Trabalho} & \textbf{Objeto e contexto} & \textbf{Achado principal} & \textbf{Limite para este problema} \\
\midrule
Sorz-Nechay et al. \cite{sorznechay2025} & Encaminhamento de pacientes pediátricos verdes e azuis a serviços externos, Viena & Redução do volume e 1,1\% de internações entre os seguidos & Depende de rede externa com capacidade; sem componente digital \\
Storm-Versloot et al. \cite{stormversloot2014} & Implantação do MTS, hospital universitário, Países Baixos & Melhor distribuição da espera, sem redução geral & Triagem isolada não cria capacidade \\
Joseph et al. \cite{joseph2023} & Revisão de 9 estudos sobre quiosques em emergência & Ganhos em etapas administrativas e pré-triagem & Não conecta o quiosque a uma alternativa de cuidado \\
Dickson et al. \cite{dickson2022} & Autotriagem eletrônica comparada à enfermagem, Reino Unido & Concordância baixa (kappa 0,14) & Desaconselha triagem pelo dispositivo \\
Gøeg et al. \cite{goeg2018} & Arquitetura modular com HL7 FHIR, Dinamarca & Núcleo de baixo acoplamento implementado e testado por unidade & Sem avaliação clínica; sem perfis brasileiros \\
Jat et al. \cite{jat2024} & Revisão de 17 arquiteturas de telemedicina & Segurança, interoperabilidade e escalabilidade recorrentes & Regulação tratada de forma superficial \\
\midrule
\textbf{Esta proposta} & Teleconsulta em totem após triagem presencial, SUS & Arquitetura rastreável com critérios clínicos e mapeamento para a RNDS & Avaliação apenas analítica \\
\bottomrule
\end{tabularx}
\fonte{Elaboração própria, a partir das fontes citadas.}
\end{table}

O que distingue a proposta é a combinação de quatro elementos que, isoladamente, já aparecem na literatura: a classificação de risco presencial como pré-requisito, o totem como canal de acesso a uma consulta e não como instrumento de triagem, critérios clínicos explícitos de elegibilidade e retorno, e a especificação da integração com os modelos e perfis nacionais.
